\begin{abstract}

It is known that neural networks have the problem of being over-confident when directly using the output label distribution to generate uncertainty measures. Existing methods mainly resolve this issue by retraining the entire model to impose the uncertainty quantification capability so that the learned model can achieve desired performance in accuracy and uncertainty prediction simultaneously. However, training the model from scratch is computationally expensive and may not be feasible in many situations. In this work, we consider a more practical post-hoc uncertainty learning setting, where a well-trained base model is given, and we focus on the uncertainty quantification task at the second stage of training. We propose a novel Bayesian meta-model to augment pre-trained models with better uncertainty quantification abilities, which is effective and computationally efficient. Our proposed method requires no additional training data and is flexible enough to quantify different uncertainties and easily adapt to different application settings, including out-of-domain data detection, misclassification detection, and trustworthy transfer learning. We demonstrate our proposed meta-model approach's flexibility and superior empirical performance on these applications over multiple representative image classification benchmarks.


% Next, we propose a simple and effective uncertainty learning method based on three crucial insights: (1) The use of multiple layers of feature representation is important for uncertainty quantification, especially for OOD tasks. (2) Leveraging Dirichlet techniques to help model different kinds of uncertainties, such as epistemic uncertainties. (3) There exists an over-fitting issue in uncertainty learning similar to supervised learning that needs to be addressed to achieve desired performance.
\end{abstract}